% GENERATED FILE. Edit canonical lessons, then run npm run generate:books.
% canonical-source-hash: fnv1a64:02bfd727dea53875
% canonical-lessons: ZH-C92-ye, ZH-W92-ye, ZH-C92-caidan, ZH-C92-fuwuyuan, ZH-C92-biede, ZH-C92-practice

\chapter{Me Too}
\label{ch:zh-me-too}
\hlchaptermodality{92}

\begin{chapteropening}
\textbf{By the end of this chapter:} \emph{I can order at a table and say 'me too'.}

Ordering at a table, the menu, the waiter and anything else, and 'me too' with \zh{也}, built on one new written character: \zh{也}.
\end{chapteropening}

\section[\texorpdfstring{\zh{也} (yě)}{也 (yě)}]{\zh{也} (yě) — also; too}

\label{lesson:ZH-C92-ye}



\begin{quote}
A new chapter, and it opens with the word you need when somebody else has ordered first.

\begin{itemize}
\raggedright
  \item \emph{Recall:} say \zh{我想要一点儿水} — \emph{wǒ xiǎngyào yìdiǎnr shuǐ}
  \item \emph{Recall:} say \zh{一点儿} — \emph{yìdiǎnr}, a little
  \item \emph{Recall:} say \zh{他是学生} — \emph{tā shì xuésheng}, he's a student
\end{itemize}
\end{quote}

\subsection*{You'll want to know: \zh{也}}
\textbf{\zh{也}} — \emph{yě} — \textquotedblleft{}also; too\textquotedblright{}.

\begin{quote}
\textbf{\zh{我也要}} \emph{wǒ yě yào} — Me too; I'd like one too.
\end{quote}

\begin{quote}
\textbf{\zh{他也是学生}} \emph{tā yě shì xuésheng} — He's a student too.
\end{quote}

The next lesson puts the character in your hand. You have written it before, as the right half of \zh{他}.

\begin{cousinweb}
\textbf{\zh{也}} (\emph{yě}) is inherited from Middle Chinese \emph{yaeX} (Guangyun, in Baxter's transcription), glossed 'a particle that closes a sentence'. 'Also' is not given there, so how it came to mean 'also' is not settled.
\end{cousinweb}

\begin{grammarlens}[title={What you've built}]
\zh{也} goes in front of the verb: \zh{我也要}, \zh{他也是学生}. English puts \emph{too} at the end; Mandarin does not.
\end{grammarlens}

\subsection*{Your turn}
\begin{itemize}
\raggedright
  \item \emph{Say it:} \emph{yě}
  \item \emph{Say it:} \emph{yě}, once more
  \item \emph{Say it:} \zh{我也要} — \emph{wǒ yě yào} — Me too.
  \item \emph{Say it:} \zh{他也是学生} — \emph{tā yě shì xuésheng} — He's a student too.
  \item \emph{Recall:} say \emph{wǒ xiǎngyào yìdiǎnr shuǐ}, then answer it for yourself with \zh{我也要}
\end{itemize}

\begin{tcolorbox}[breakable,colback=teal!4,colframe=teal!35!black,title={Before you move on}]
What does \zh{也} mean? (\textquotedblleft{}Also; too\textquotedblright{}.) Where does it go? (\textbf{In front of the verb: \zh{我也要}}.)
\end{tcolorbox}
\section[\texorpdfstring{\zh{也} (yě)}{也 (yě)}]{\zh{也} — the right half of \zh{他}}

\label{lesson:ZH-W92-ye}



\begin{quote}
The word first, then the character inside it.

\begin{itemize}
\raggedright
  \item \emph{Recall:} say \emph{yě} — \textbf{R1}, one lesson back, before it decays
\end{itemize}

\begin{quote}
\zh{他} \zh{他们}
\end{quote}

\emph{Read it:} each one aloud, with nothing said first — \emph{tā}, \emph{tāmen}
\end{quote}

\subsection*{Script — the right half of \zh{他}}
\hlblockfigure{\detokenize{figures/ZH-W92-ye-filmstrip.pdf}}{How \zh{也} is written, stroke by stroke}

\begin{quote}
\zh{也}
\end{quote}

\textbf{One piece, and your hand already knows it: the right half of \zh{他}, stroke for stroke.}

\textbf{A hooked turn}: a horizontal that rises a little, turns down the right side and hooks to the left.

\textbf{A short vertical} from above it, down through it.

\textbf{The long bend}: down the left from above the horizontal, along the foot to the right, and a hook up at the end.

Three strokes.

\subsection*{Your turn}
\begin{itemize}
\raggedright
  \item \emph{Write it:} \zh{也}, all three, with the model in front of you
  \item \emph{Cover:} write it again from memory, and count the pieces, not the strokes
\end{itemize}

\begin{tcolorbox}[breakable,colback=teal!4,colframe=teal!35!black,title={Before you move on}]
How many strokes in \zh{也}? (\textbf{Three}.) Where have you written them before? (\textbf{On the right of \zh{他}}.) Which stroke ends with a hook up? (\textbf{The last, the long bend}.)

Source: \href{https://raw.githubusercontent.com/chanind/hanzi-writer-data/68d10a4b21150cae5e1ebbd223eed289cf32d90c/data/\%E4\%B9\%9F.json}{Hanzi Writer Data \zh{也}, pinned stroke paths and medians}.
\end{tcolorbox}
\section[càidān]{càidān — menu}

\label{lesson:ZH-C92-caidan}



\begin{quote}
Before the new one: read the last lesson's character, then two words for food.

\begin{itemize}
\raggedright
  \item \emph{Read it:} \zh{我也要}, aloud — \emph{wǒ yě yào}
  \item \emph{Recall:} say \emph{cài}, a dish — \textbf{R4}, recognition at distance
  \item \emph{Recall:} say \emph{diǎncài}, to order food — \textbf{R3}, durable
\end{itemize}
\end{quote}

\subsection*{You'll want to know: càidān}
\emph{càidān} — \textquotedblleft{}menu\textquotedblright{}.

\begin{quote}
\emph{qǐng gěi wǒ càidān} — The menu, please.
\end{quote}

Its characters come later; for now it is a word for the ear.

\begin{cousinweb}
\emph{càidān} is built from cài 'a dish', the \emph{cài} you know, and dān, both inherited from Middle Chinese, \emph{tshojH} and \emph{tan} (Guangyun, in Baxter's transcription), glossed 'greens fit to eat' and 'single'. A sheet that lists the dishes is the modern sense of the pair.
\end{cousinweb}

\begin{grammarlens}[title={What you've built}]
\emph{qǐng gěi wǒ càidān}: \zh{请}, then \zh{给} + the person + the thing, the pattern of the chapter on giving.
\end{grammarlens}

\subsection*{Your turn}
\begin{itemize}
\raggedright
  \item \emph{Say it:} \emph{càidān}
  \item \emph{Say it:} \emph{càidān}, once more
  \item \emph{Say it:} \emph{qǐng gěi wǒ càidān} — The menu, please.
  \item \emph{Recall:} say \emph{diǎncài}, then \emph{càidān} again
\end{itemize}

\begin{tcolorbox}[breakable,colback=teal!4,colframe=teal!35!black,title={Before you move on}]
What does \emph{càidān} mean? (\textquotedblleft{}Menu\textquotedblright{}.) Ask for it. (\emph{qǐng gěi wǒ càidān}.)
\end{tcolorbox}
\section[fúwùyuán]{fúwùyuán — waiter; attendant}

\label{lesson:ZH-C92-fuwuyuan}



\begin{quote}
Before the new one: last lesson's word, then a place to eat, then a guest.

\begin{itemize}
\raggedright
  \item \emph{Recall:} say \emph{càidān}, the menu — \textbf{R1}, one lesson back
  \item \emph{Recall:} say \emph{fàndiàn}, a restaurant — \textbf{R4}, recognition at distance
  \item \emph{Recall:} say \zh{客人} — \emph{kèrén}, a guest
\end{itemize}
\end{quote}

\subsection*{You'll want to know: fúwùyuán}
\emph{fúwùyuán} — \textquotedblleft{}waiter; attendant\textquotedblright{}.

\begin{quote}
\emph{fúwùyuán, wǒ yào càidān} — Excuse me, the menu, please.
\end{quote}

Its characters come later; for now it is a word for the ear.

\begin{cousinweb}
\emph{fúwùyuán} is built from fú 'to serve', wù 'a task' and yuán, all inherited from Middle Chinese, \emph{bjuwk}, \emph{mjuH} and \emph{hjwen} (Guangyun, in Baxter's transcription), glossed 'to serve', 'affairs' and 'a number of things'. 'A member of staff' is a later sense of yuán.
\end{cousinweb}

\begin{grammarlens}[title={What you've built}]
In a restaurant or a shop, \emph{fúwùyuán} is also how you call the person serving: say it, then what you want.
\end{grammarlens}

\subsection*{Your turn}
\begin{itemize}
\raggedright
  \item \emph{Say it:} \emph{fúwùyuán}
  \item \emph{Say it:} \emph{fúwùyuán}, once more
  \item \emph{Say it:} \emph{fúwùyuán, wǒ yào càidān} — Excuse me, the menu, please.
  \item \emph{Recall:} say \emph{fàndiàn}, then \emph{fúwùyuán} again
\end{itemize}

\begin{tcolorbox}[breakable,colback=teal!4,colframe=teal!35!black,title={Before you move on}]
What does \emph{fúwùyuán} mean? (\textquotedblleft{}Waiter; attendant\textquotedblright{}.) Call one and ask for the menu. (\emph{fúwùyuán, wǒ yào càidān}.)
\end{tcolorbox}
\section[biéde]{biéde — something else; other}

\label{lesson:ZH-C92-biede}



\begin{quote}
Before the new one: last lesson's word, then the particle at the end of this one.

\begin{itemize}
\raggedright
  \item \emph{Recall:} say \emph{fúwùyuán}, a waiter — \textbf{R1}, one lesson back
  \item \emph{Recall:} say \emph{dōngxi}, things — \textbf{R4}, recognition at distance
  \item \emph{Recall:} say \zh{我的书} — \emph{wǒ de shū}, my book
\end{itemize}
\end{quote}

\subsection*{You'll want to know: biéde}
\emph{biéde} — \textquotedblleft{}something else; other\textquotedblright{}.

\begin{quote}
\emph{nǐ yào biéde ma} — Would you like anything else?
\end{quote}

\begin{quote}
\emph{wǒ xiǎngyào biéde dōngxi} — I'd like something else.
\end{quote}

Its first character comes later; for now it is a word for the ear.

\begin{cousinweb}
\emph{biéde} is built from bié 'other; apart' and de, the \zh{的} you write. bié is inherited from Middle Chinese \emph{bjet} (Guangyun, in Baxter's transcription), glossed 'different; to part'; where de comes from is not settled.
\end{cousinweb}

\begin{grammarlens}[title={What you've built}]
\emph{nǐ yào biéde ma} is the waiter's question; \emph{bú yào, xièxie} is the answer when you have all you need.
\end{grammarlens}

\subsection*{Your turn}
\begin{itemize}
\raggedright
  \item \emph{Say it:} \emph{biéde}
  \item \emph{Say it:} \emph{biéde}, once more
  \item \emph{Say it:} \emph{nǐ yào biéde ma} — Would you like anything else?
  \item \emph{Say it:} \zh{不要} \zh{谢谢} — \emph{bú yào, xièxie} — No, thank you.
  \item \emph{Recall:} say \emph{dōngxi}, then \emph{biéde} again
\end{itemize}

\begin{tcolorbox}[breakable,colback=teal!4,colframe=teal!35!black,title={Before you move on}]
What does \emph{biéde} mean? (\textquotedblleft{}Something else; other\textquotedblright{}.) Answer the waiter's question politely. (\textbf{\zh{不要} \zh{谢谢}}.)
\end{tcolorbox}
\section[Practice]{\zh{我也要}}

\label{lesson:ZH-C92-practice}



\begin{quote}
The chapter closes here, at a table for two.

\begin{itemize}
\raggedright
  \item \emph{Recall:} say \emph{biéde}, something else — \textbf{R1}, one lesson back
  \item \emph{Recall:} say \emph{jiǎozi}, dumplings — \textbf{R3}, durable
  \item \emph{Recall:} say \emph{xūyào}, to need — \textbf{R2}
\end{itemize}
\end{quote}

\subsection*{The exchange}
\begin{quote}
\textbf{A:} \emph{fúwùyuán, qǐng gěi wǒmen càidān} — Excuse me, could we have the menu?
\end{quote}

\begin{quote}
\textbf{B:} \zh{好} \zh{你们喝什么}
\end{quote}

\begin{quote}
\emph{hǎo, nǐmen hē shénme} — Of course. What will you drink?
\end{quote}

\begin{quote}
\textbf{A:} \zh{我要水}
\end{quote}

\begin{quote}
\emph{wǒ yào shuǐ} — Water for me.
\end{quote}

\begin{quote}
\textbf{C:} \zh{我也要}
\end{quote}

\begin{quote}
\emph{wǒ yě yào} — Me too.
\end{quote}

\begin{quote}
\textbf{B:} \emph{nǐmen yào biéde ma} — Anything else?
\end{quote}

\begin{quote}
\textbf{A:} \emph{wǒ yào jiǎozi} — Dumplings for me.
\end{quote}

\begin{quote}
\textbf{C:} \emph{wǒ yě yào jiǎozi} — Dumplings for me too.
\end{quote}

\begin{quote}
\textbf{B:} \zh{好}
\end{quote}

\begin{quote}
\emph{hǎo} — Right you are.
\end{quote}

Ordering is \emph{yào} and the thing; agreeing is \zh{也} in front of it. Nobody says \emph{I} would like: \zh{我要} does the job.

\begin{tcolorbox}[breakable,skin=enhanced,colback=orange!6,colframe=orange!45!black,boxrule=0.5pt,arc=1mm,left=6pt,right=6pt,top=4pt,bottom=4pt,fonttitle=\bfseries,title={Read this}]
\emph{Read it:} the exchange through once without stopping, then answer the question below
\end{tcolorbox}

\begin{grammarlens}[title={What you've built — ordering at a table}]
\begin{itemize}
\raggedright
  \item \emph{qǐng gěi wǒmen càidān} — the menu, please.
  \item \textbf{\zh{我要水}} \emph{wǒ yào shuǐ} — water for me.
  \item \textbf{\zh{我也要}} \emph{wǒ yě yào} — me too.
  \item \emph{nǐmen yào biéde ma} — anything else?
\end{itemize}

\textbf{Listen.} \emph{Nǐmen hē shénme?}

\textbf{Speak.} Order a drink, then agree with somebody else's order using \zh{也}.

\emph{Write it:} \zh{也} — the hooked turn, the short vertical, the long bend
\end{grammarlens}

\subsection*{Your turn}
\begin{itemize}
\raggedright
  \item \emph{Say it:} \zh{我也要} — \emph{wǒ yě yào} — Me too.
  \item \emph{Say it:} \zh{请给我们水} — \emph{qǐng gěi wǒmen shuǐ} — Water for us, please.
  \item \emph{Say it:} the whole exchange, every part
  \item \emph{Write it:} \zh{我也要}
\end{itemize}

\begin{tcolorbox}[breakable,colback=teal!4,colframe=teal!35!black,title={Before you move on}]
Say me too. (\textbf{\zh{我也要}}.) Ask for the menu. (\emph{qǐng gěi wǒ càidān}.) Call the waiter. (\emph{fúwùyuán}.) Where does \zh{也} go? (\textbf{In front of the verb}.)
\end{tcolorbox}
